\section{Introduction}
\label{section:introduction}

This primer aims to introduce social science readers to an exciting literature exploring how deep neural networks can be used to estimate causal effects. In recent years, both causal inference frameworks and deep learning have seen rapid adoption across science, industry, and medicine. Causal inference has a long tradition in the social sciences, and social scientists are increasingly exploring the use of machine learning for causal inference \citep{Athey2016,wager2018, chernozhukov2016double}. Nevertheless, deep learning remains conspicuously underutilized by social scientists compared to other ML approaches, both for causal inference and more generally.

The deep learning revolution has been spurred by the flexibility and expressiveness of these models. Neural networks are nearly non-parametric and can theoretically approximate any continuous function \citep{Cybenko1989ApproximationBS}, making them well suited for both classification and regression tasks. Furthermore, they can be configured with different architectures and objectives to learn from a variety of quantitative data as well as text, images, video, networks, and speech. These advantages allow them to learn vector ``representations" of complex data with emergent properties. Simple examples of representation learning include the Word2Vec algorithm that discovers semantic relationship between words in texts, or face classification models that learn vectors describing facial features \citep{mikolov2013}. More recently, generative models like DALL-E, Stable Diffusion, and ChatGPT have shown how coherent text passages and life-like images can be reconstructed from learned representations. 

Here we explore the potential for leveraging these advantages to estimate causal effects. Causal inference frameworks are non-parametric, but the linear models traditionally used to estimate causal effects require strong parametric assumptions. In contrast, the nearly non-parametric nature of neural networks allows us to estimate smooth response surfaces that capture heterogeneous treatment effects for individual units with low bias.\footnote{Neural networks can have hundreds to billions of parameters making them effectively non-parametric. The risks of overparameterization of neural networks are discussed in Section \ref{section:primerdl}.} The ability of these models to learn from complex data means we can extend causal inference to new settings where confounding is complicated, time-varying, or even encoded in texts, graphs, or images (see Box \ref{Box:nontradex} for hypothetical examples). Lastly given the right objectives, neural networks promise to learn deconfounded representations of data, presenting a new strategy for treatment modeling.

This primer synthesizes existing literature on deep causal estimators, but it is not a review; its goals are fundamentally pedagogical and prospective rather than retrospective. In Section \ref{section:primerdl}, we introduce social scientists to the fundamental concepts of deep learning, and the basic workflow for building and training their own deep neural networks within a supervised learning framework. For readers unfamiliar with causal inference, Section \ref{section:CI} introduces the assumptions of causal identification and three fundamental estimation strategies within the selection on observables design: matching, outcome modeling, and inverse propensity score weighting. Machine learning models often perform poorly in both theory and practice when only one of these strategies is employed, so we also introduce the concept of double robustness.

Section \ref{section:mainbody} is the main body of the article. Here we introduce four related deep learning models for the estimation of heterogeneous treatment effects: the S-learner, T-learner, TARNet and Dragonnet \citep{Shalit2017,Shi2019}. Although this literature is rapidly evolving, these four models are sufficient to illustrate how traditional estimation strategies can be used in creative ways that leverage the key strengths of neural networks (i.e., deconfounding through representation learning, semi-parametric inference). Section \ref{section:interpretation} deals with the practical considerations of building confidence intervals and interpreting neural networks. These guidelines are concretized in the \href{https://github.com/kochbj/Deep-Learning-for-Causal-Inference}{companion online tutorials}, which show readers how to implement and interpret the models described in Section \ref{section:mainbody} in Tensorflow 2 and PyTorch.

In Section \ref{section:nontraditional}, we focus on the future of deep causal information: estimators that can disentange counfounding relationships embedded within texts, images, graphs, or time-varying data. In the interest of clarity, we give hypothetical examples of the types of questions social scientists might answer with these models, and briefly describe ongoing research on each of these modalities. For fuller treatments of some of these models, see the appendix. We conclude with a discussion of how neural networks fit into the broader literature on machine learning for causal inference (Section \ref{section:conclusions}).

The primer makes multiple contributions. First, it is one of the first pieces in the sociological literature to introduce the fundamentals of deep learning not only at a conceptual level (e.g, backpropagation, representation learning), but at a practical one (e.g., validation, hyperparameter tuning). Our recommendations for training and interpreting neural networks are supported by heavily annotated tutorials that teach readers without prior familiarity with deep learning how to build their own custom models in Tensorflow 2 and PyTorch. Second, we use this foundation and select examples to build intuition on how the core strengths of deep learning can be leveraged for causal inference. Finally, we highlight future directions for this literature and argue why the future of causal estimation runs through deep learning.

\begin{TCBBox}[label=Box:nontradex]{Example Scenarios for Causal Inference with Non-Traditional Data}
\textbf{Text.} As a motivating example, \citet{Veitch2019UsingTE} consider the effect of the author's reported gender ($T$) on the number of upvotes a Reddit post receives ($Y$). However gender may also ``affect the text of the post, e.g., through tone, style, or topic
choices, which also affects its score [($X$)]." Controlling for a representation of the text would allow the analyst to more accurately estimate the direct effect of gender.

\textbf{Images.} \citet{todorov2005inferences} showed that split second-judgments of a politician's competence ($T$) from pictures ($X$) of their face is predictive of their electability ($Y$). When attempting to replicate this study using machine learning classifiers rather than human classifiers, \citet{joo2015automated} suggest that the age of the face ($Z$) is a not-so-obvious confounder: while older individuals are more likely to appear competent, they are also more likely to be incumbents. Even if age is unknown, using neural networks to control for confounders implicitly encoded in the image (like age) could reduce bias.

\textbf{Networks.} \citet{nagpal2020} explore the question of which types of prescription opioids (e.g., natural, semi-synthetic, synthetic) ($T$) are most likely to cause long term addiction ($Y$). Because of predisposition to different injuries, type of employment ($X$) could be a common cause of both treatment and outcome. Suppose job type is unobserved, but we know that patients are likely to associate with coworkers through homophily. To capture some of the effects of this latent unobserved confounder, analysts might choose to control for a representation of the patient's position in their social network when estimating the causal effect.

\end{TCBBox}